\section{The Harness Plane Advances}
\sectionkicker{AGENT HARNESS}
\begin{wideflow}
{\Large\bfseries Agents APIとFusion、実行基盤の2つの動き}\par
\smallskip
{\color{SurveyMuted}9月10日のAgents APIパブリックベータと、9月11日17時（世界時）に通常ウィンドウ内と確認したFusionを、提供形態の区別とともに並べて読む。}\par
\end{wideflow}

9月10日、OpenAIはAgents APIをパブリックベータとして公開した\cite{agentsapi}。Codexのエージェントハーネスと実行基盤を開発者向けに束ねる内容で、タスク・モデル・ツール・実行環境を指定する1回のAPI呼び出しで本番利用可能なエージェントを起動できるとされる。例示ではGPT-6 Astraが設定の一例として示されるが、API自体をAstraのAPIとは呼ばない。実行環境はOpenAIが管理するサンドボックス、自社インフラ、Blaxel・Cloudflare・Daytona・DigitalOcean・E2B・Modal・Oracle・Runloop・Vercelといったパートナーから選べるとされる。長時間セッション向けのコンテキスト自動圧縮、必要に応じたツール定義の読み込みや並列・連結実行、独立したコンテキストを持つ並列サブエージェント、モデルリリースに合わせたバージョン管理、オープンソースのCodexハーネスが土台として示される。パブリックベータの適用範囲や正式提供時期、レート制限・提供地域は資料に見当たらない。連携の深さやVPC・セキュリティの詳細はベンダー側の説明にとどめ、顧客の言葉も選ばれた紹介として扱う。API利用の追加料金はなく、トークンとツールの利用分を支払う構成とされる。

9月11日のCognitionのFusionは、Devin DesktopとCLIで動く並列型のエージェントハーネスである\cite{fusion}。リードエージェントが計画とレビューを担い、SWE-2が実行役のサイドキックに回る組み合わせで、各エージェントが独自のコンテキストを持って並走する。推奨構成はFable 5.1とSWE-2の組み合わせとされる。公開時刻は一次ページのメタデータ2系統がともに9月11日10時（米太平洋夏時間）を示し、世界時では9月11日17時となる。通常ウィンドウの終了が同日22時（世界時）のため、5時間の余裕で通常ウィンドウに収まる。日付のみの推定は退け、この時刻を根拠とする。効率見出しの39\%は評価表の中の最大値であり、一律の改善ではない。DeepSWEやTerminal-Bench、SWE-Atlas Q\&A、Vals移行、FrontierCode拡張で、スコアをほぼ保ったままコストを抑えるベンチマークごとの差分が示されるが、評価手法はベンダー側の説明の範囲で扱う。検証していないリード・サイドキックの組み合わせへの一般化は成立していない。

\begin{claimboundary}[提供形態と効率の境界]
正式提供時期やレート制限・提供地域は資料にない。39\%は最大値であり一律ではない。評価手法はベンダー側のまとめの範囲である。
\end{claimboundary}

\begin{communitynote}[X / Community observation --- context only]
9月11日16時16分（世界時）の公式発信に、締め切り前の独立した反応が重なる様子が観察される\cite{w37x-coding-official,w37x-coding-ind}。うかがえるのは音声対応の開発環境への関心であり、仕様やベンチマークの根拠にはならない。
\end{communitynote}
